\documentclass[11pt, a4paper]{article}
\input{../Config.tex}
\usepackage{hyperref}
\renewcommand{\thesubsection}{}
\renewcommand{\thesection}{}

\begin{document}
\pagestyle{fancy}
% 
\lhead{Bernat Esteve Sagarra}
\rhead{Problemes de teoria de nombres}
% ------------------------------------------------------
% -------------------- Exercici 1 ----------------------
% ------------------------------------------------------
    \begin{tcolorbox}[title=\section{FP1 Exercici 1}]
        L'anell dels \textit{enters de Gauss} es defineix com
        \[
            \ZZ[i]=\{a+ib\mid a,b\in \ZZ\},
        \]
        on $i=\sqrt{-1}$. Anàlogament, anomenarem
        \[
            \QQ(i)=\{a+ib\mid a,b\in \QQ\}
        \]
        el cos dels \textit{nombres racionals de Gauss}. Definim la \textit{norma} d'un nombre de Gauss com 
        \[
            N(\alpha)=\alpha\bar{\alpha}=a^2+b^2.
        \]
        \begin{enumerate}
            \item Demostreu que $\ZZ[i]$ és un domini euclidià respecte a la norma $N:\ZZ[i]\to\ZZ_{\geq0}$.
            \item Demostreu que $N(\alpha\beta)=N(\alpha)N(\beta)$ per a tot $\alpha,\beta\in \ZZ[i]$.
            \item Demostreu que $\alpha\in\ZZ[i]$ és unitat si i només si $N(\alpha)=1$.
            \item Sigui $p$ un primer senar de $\ZZ$. Demostreu que $p$ no és un element primer de $\ZZ[i]$ si i només si $p\equiv1\pmod 4$.
            \item (Fermat 1640) Sigui $p$ un primer senar de $\ZZ$. Demostreu que $p$ es pot expressar com la suma dels quadrats de dos nombres enters si i només si $p\equiv1\pmod4$.
            \item Demostreu que llevat d'associats, els elements primers de $\ZZ[i]$ són:
            \begin{enumerate}
                \item $1+i$,
                \item $a+ib$, on $a,b\in\ZZ$ són tals que $a^2+b^2$ és un primer $p\equiv1\pmod4$ de $\ZZ$, i $a>|b|>0$.
                \item $p$. on $p\equiv3\pmod 4$ és un primers de $\ZZ$.
            \end{enumerate}
        \end{enumerate}
    \end{tcolorbox}
    \subsection{Exercici 1.1}
    Volem veure que si tenim $a,b\in\ZZ[i]$ ($b\neq0$), aleshores podem trobar $q,r\in \ZZ[i]$ tals que $a=bq+r$ amb $N(r)<N(b)$. Per construir tals $q,r$, prenem $c+di=a/b\in\QQ[i]$. Siguin $c'$ i $d'$ els enters més pròxims a $c$ i $d$ respectivament (és a dir, $c-c'=\eps_c$ i $d-d'=\eps_d$ tals que $|\eps_c|, |\eps_d|\leq\half$). Sigui ara $q=c'+d'i$ i $r=b\eps_c+b\eps_di$
    \[
        a=b(c+di)=b(c'+d'i)+ (b\eps_c+b\eps_di)=bq+r.
    \]
    Ara si veiem que $N(r)<N(b)$ ja estarem, però:
    \[
        N(r)=N(b\eps_c+b\eps_di)\overset{\text{Apartat }2}{=}N(b)(\eps_c^2+\eps_d^2)\leq N(b)/2< N(b)
    \]
    Que és el que volíem veure.
    \subsection{Exercici 1.2}
    Si $\alpha=a+bi$ i $\beta=c+di$, aleshores 
    \begin{align*}
        N(\alpha\beta)&=N(ac-bd+bci+adi)=\\
        &=(ac-bd)^2+(bc+ad)^2=\\
        &=a^2c^2-2abcd+b^2d^2+b^2c^2+2abcd+a^2d^2=\\
        &=(a^2+b^2)(c^2+d^2)=\\
        &=N(\alpha)N(\beta)
    \end{align*}
    \subsection{Exercici 1.3}
    Per veure la primera implicació ($\impliedby$), notem que si $\alpha$ és unitat, aleshores $N(\alpha\alpha^{-1})=N(1)=1$, però $N(\alpha)N(\alpha^{-1})$=1, que ens diu que ambdues són iguals a 1.
    % notem que $N$ és un morfisme de grups que envia $(\ZZ[i],\times)\to(\ZZ_{\geq0},\times)$, per tant les unitats de $\ZZ[i]$ s'envien a les unitats de $\ZZ_{\geq0}$.\par
    Per veure la segona implicació, si $N(\alpha)=1$ aleshores per la definició de norma $\alpha\bar{\alpha}=1$, per tant, hem trobat l'element invers de $\alpha$.
    \subsection{Exercici 1.4}
    Sigui $p$ un primer senar de $\ZZ$ però que no és primer a $\ZZ[i]$, aleshores tindríem $p=ab$, i fent normes als 2 costats aconseguim: $p^2=N(a)N(b)$. Com que $p$ és primer a $\ZZ$, tenim 2 possibilitats: $N(a)=N(b)=p$, o $N(a)=p^2$ i $N(b)=1$. En el segon cas tenim que $b$ és unitat, i per tant, no ens interessa. Suposem que $N(a)=N(x+iy)=x^2+y^2=p$ i mirant mòdul 4 tenim $p=x^2+y^2\equiv 0,1,2\pmod4$, però com que $p$ és senar, $p\equiv1\pmod4$.\par

    Suposem que $p\equiv1\pmod4$, aleshores sabem que $\left(\frac{-1}{p}\right)=1$, per tant, existeix alguna solució a l'equació $x^2+1\equiv0\pmod p$. Sigui $x\in\ZZ$ un enter que satisfaci aquesta equació, aleshores $p\mid x^2+1=(x+i)(x-i)$. Suposem que $p$ és primer, aleshores $p\mid x+i$ o $p\mid x-i$, per tant $p(a+bi)=x\pm i$, i fixant-nos en la part imaginària, obtenim que $pb=\pm1$, que no pot ser.
    \subsection{Exercici 1.5}
    Si $p$ es pot expressar com la suma de 2 quadrats (pel mateix argument d'abans de mirar-te les coses mòdul 4), obtens que $p\equiv 1\pmod4$.\par
    Si $p$ és congru amb 1 mòdul 4, aleshores $p=\alpha\beta$, amb $\alpha, \beta\in \ZZ[i]$ (cap dels dos és una unitat), i per tant $N(\alpha)=p$, és a dir, $N(x+yi)=x^2+y^2=p$, que és el que volíem veure.
    \subsection{Exercici 1.6}
    Recordem que si $N(x)$ és un irreductible, aleshores $x$ és irreductible.
    \begin{dem}[Si $N(x)$ és irreductible, aleshores $x$ és irreductible.]
        Suposem que $x=ab$, mirant normes ens queda que $N(x)=N(a)N(b)$, però $N(x)$ és irreductible, és a dir que $N(a)=1$ o $N(b)=1$, i per tant, almenys un de $a$ i $b$ són unitats.
    \end{dem}
    A partir d'aquest mini lema, podem veure ràpidament que els primers de la forma $1+i$ i $a+ib$ que satisfan $N(a+ib)=p$ són primers (com que estem en un anell de factorització única, primer = irreductible). A més, ja hem vist que els primers de $\ZZ$ de la forma $p\equiv3\pmod4$ són primers de $\ZZ$.\par
    Ara hem de veure que no n'hi ha més. Aleshores si veiem que les úniques normes possibles per un primer $z$ són $N(z)=p$ o $N(z)=p^2$, ja estarem, ja que $N(z)=p$ sempre és primer, i si $N(z)=p^2$, aleshores ha de ser forçosament $z=p\equiv3\pmod4$.\par
    Sigui $z=a+bi\in \ZZ[i]$ primer. Considerem $N(z)$: podria ser que $N(z)$ ja sigui un nombre primer, i en aquest cas ja estaríem; en cas contrari, sigui $p\mid N(z)$. Per estudiar aquest cas, considerarem primer $p\equiv3\pmod 4$, i després $p\equiv 1\pmod 4$.\par
    Suposem que $p\equiv 3\pmod4$, aleshores $p\mid z\ol z$, per tant (com que $p$ és primer a $\ZZ[i]$), tenim que $p\mid z$ o $p\mid \ol z$, però com que $p=\ol p$, tenim que de fet passen els 2 casos, i en particular $p\mid z$, que ens diu que $z\sim p$ (és un associat de $p$).\\
    El segon cas és si \underline{$p\equiv 1\pmod 4$}, és a dir, tenim que $p\mid z\ol z$, però com que $p\equiv 1\pmod 4$, tenim que $p=\pi\ol\pi$ amb $\pi$ primer de $\ZZ[i]$; i com que $\pi\mid z\ol z$, tenim $\pi\mid z$ o $\pi\mid \ol z\iff \ol\pi\mid z$, i per tant, $z\sim \pi$ o $z\sim \ol\pi$; que és el que volíem veure.
    \newpage
% ------------------------------------------------------
% -------------------- Exercici 2 ----------------------
% ------------------------------------------------------
    \begin{tcolorbox}[title=\section{FP1 Exercici 2}]
        Diem que $(x,y,z)\in \ZZ_{>0}^3$ és una \textit{terna pitagòrica} si $x^2+y^2=z^2$.
        \begin{enumerate}
            \item Demostreu que per $\alpha\in\QQ(i)$:
            \[
                N(\alpha)=1\iff \alpha=\frac{\beta}{\ol\beta}\text{ per algun }\beta\in\QQ(i)^\times
            \]
            \item Deduïu de l'apartat anterior que si $(x,y,z)$ és una terna pitagòrica, aleshores existeix un enter positiu $d$ i dos enters coprimers $u$ i $v$ tals que
            \[
            x=d(v^2-u^2),\quad y=2duv,\quad z=d(u^2+v^2),
            \]
            llevat de permutació de $x$ i $y$.
            \item Doneu una demostració geomètrica de l'apartat anterior seguint els següents passos. Considereu $C:X^2+Y^2-1$ i $L_m:Y-m(X+1)$ per $m\in\QQ$. Demostreu que existeix una bijecció entre els punts de coordenades racionals de $C\setminus \{(-1,0)\}$ amb les ternes pitagòriques primitives.
        \end{enumerate}
    \end{tcolorbox}
    \subsection{Exercici 2.1}
    Notem que això és exactament un cas particular del \textit{teorema 90 de Hilbert}. Adaptaré la demostració que hi ha als apunts del Xevi, en el nostre cas particular.
    La implicació $\impliedby$ és evident. Considerem ara $\alpha\in \ZZ[i]$ tal que $N(\alpha)=1$. Si $\alpha=-1$, aleshores $\alpha =\frac{i}{-i}$. Suposem ara que $\alpha\neq -1$, sigui $\beta = 1+\alpha$, aleshores:
    \[
        \alpha\ol\beta=\alpha+ \alpha\ol\alpha\overset{\alpha\ol\alpha=1}{=}1+\alpha=\beta
    \]
    I com que $\beta\neq0$, tenim que $\alpha=\frac{1+\alpha}{\ol{1+\alpha}}$, que és el que volíem.
    \subsection{Exercici 2.2}
    Abans de començar, notem que reduint l'equació mòdul 2, almenys un dels termes ha de ser parell: suposem que $y$ és un nombre parell. Ara, per aplicar el que hem vist abans: simplificant $x^2+y^2=z^2$ obtenim $(x/z)^2+(y/z)^2=1$, sigui $a=x/z$ i $b=y/z$. Per tant, tenim un element $a+bi$ de norma 1, que ens diu que hi ha una element $\beta$ tal que $a+bi=\beta/\ol\beta$. Com que $y\neq0$, podem utilitzar la $\beta$ calculada a l'exercici anterior:
    \begin{align*}
        \frac{x}{z}+\frac{y}{z}i&=\frac{z+x+yi}{z+x-yi}=\\
        &=\frac{2yi}{x+z+yi}=\\
        &=2\left(\frac{x+z}{yi}+1\right)\inv
    \end{align*}
    Com que no sembla que es pugui simplificar gaire més, sigui la fracció irreductible $\frac{u}{v}=\frac{x+z}{y}$, aleshores:
    \[
        x+z=uD\qquad y=vD
    \]
    Mirant-nos l'equació original, tenim que si $y$ és un nombre parell, aleshores $x$ és parell si i només si $z$ és parell: per tant, tant $x+z$ com $y$ són parells, és a dir, $D=2D'$: $x+z = 2uD'$, tenim que $x=uD'-e$ i $z=uD'+e$. Posant això a l'equació original tenim:
    \begin{align*}
        (uD'-e)^2+(2vD')^2&=(uD'+e)^2\\
        u^2D'^2-2euD'+e^2+4v^2D'^2&=u^2D'^2+2euD'+e^2\\
        v^2D'&=eu
    \end{align*}
    Però com que hem suposat que $(u,v)=1$ tenim que $v^2\mid eu\implies v^2\mid e$; a més, $u\mid v^2D'\implies u\mid D'$, és a dir, que podem escriure $e=v^2d$, $D'=ud$ (és la mateixa constant ja que $v^2D = eu$), i per tant, tenim
    \[
        \begin{cases}
            x = u^2d-v^2d\\
            y = 2uvd\\
            z = u^2d+v^2d
        \end{cases} 
    \]
    Que és el que volíem veure.
    \subsection{Exercici 2.3}
    Sigui $P=(p/r,q/r)\neq (-1,0)$ un punt de coordenades racionals de $C$, aleshores considerem la recta que passa pel punt $P$ i pel punt $V=(-1,0)$:
    \[
        L_P:y-\frac{q}{p+r}(x+1)=0
    \]
    Notem que aquesta recta ...
    \newpage
% ------------------------------------------------------
% -------------------- Exercici 3 ----------------------
% ------------------------------------------------------
    \begin{tcolorbox}[title=\section{FP 1 Exercici 3}]
        Equació de Fermat per $n=3$. Definim l'anell dels \textit{enters d'Eisenstein} com
        \[
            \ZZ[\rho]=\{a+b\rho|a,b\in\ZZ\},\qquad\text{on }\rho=\frac{1+\sqrt{-3}}{2}
        \]
        L'objectiu d'aquest exercici és demostrar que l'equació
        \[
        x^3+y^3=z^3
        \]
        no té cap solució amb $x,y,z\in\ZZ$ tals que $xyz\neq0$. Més en general, demostrarem que, per a tot $u\in\ZZ[\rho]^\times$, l'equació
        \begin{equation}\label{FP1P3 equació de Fermat}
            x^3+y^3=uz^3
        \end{equation}
        no té cap solució amb $x,y,z\in \ZZ[\rho]$, tals que $xyz\neq0$. Podeu seguir els següents passos:
        \begin{enumerate}
            \item Demostreu que $\ZZ[\rho]$ és un anell euclidià i que $\ZZ[\rho]^\times=\langle \rho\rangle$.
            \item Proveu que $\lambda=1+\rho$ és un element primer i que $3=\rho\inv\lambda^2$. Donat $\alpha\in\ZZ[\rho]$ no nul, denotarem per $\ord_\lambda(\alpha)$ la màxima potència de $\lambda$ dividint $\alpha$. Demostreu que donats $\alpha,\beta\in\ZZ[\rho]$ se satisfà
            \[
                \ord_\lambda(\alpha+\beta)\geq\min\{\ord_\lambda(\alpha),\ord_\lambda(\beta)\}.
            \]
            Mostreu que en la desigualtat de dalt, es té igualtat si $\ord_\lambda(\alpha)\neq\ord_\lambda(\beta)$.
            \item Demostreu que l'aplicació
            \[
                \Phi:\ZZ[\rho]\to\ZZ/3\ZZ,\qquad \Phi(a+b\rho)=a-b
            \]
            és un morfisme d'anells i que indueix un isomorfisme $\ZZ[\rho]/(\lambda)\simeq\ZZ/3\ZZ$.
            \item Proveu que no existeixen solucions tals que $\lambda\nmid xyz$\\
            \textit{Indicacions: Vegeu que si $x\in\ZZ[\rho]$ i $x\equiv \pm1\pmod\lambda$, aleshores $x^3\equiv \pm 1\pmod {\lambda^4}$.}
            \item Demostreu que tota solució amb $\lambda\nmid xy$ compleix $\lambda^2\mid z$.
            \item Demostreu que, per a tota solució amb $\lambda\nmid xy$, d'entre els nombres
            \[
                \ord_\lambda(x+y),
                \ord_\lambda(\rho^2x+\rho^4y),
                \ord_\lambda(\rho^4x+\rho^2y)
            \]
            n'hi ha exactament dos que valen 1.\\
            \textit{Indicacions: Considereu la descomposició
            \[
                (x+y)(\rho^2x+\rho^4y)(\rho^4x+\rho^2y)=x^3+y^3=uz^3
            \]}
            \item Diem que una solució de (\ref{FP1P3 equació de Fermat}) és primitiva si $\gcd(x,y,z)=1$. Proveu que si $(x,y,z)$ és una solució primitiva amb $\lambda\nmid xy$, aleshores n'hi ha una altra $(x_1,y_1,z_1)$ amb $\lambda\nmid x_1y_1$ i amb $\ord_\lambda z_1<\ord_\lambda z$.
            \item Demostreu que (\ref{FP1P3 equació de Fermat}) no té cap solució amb $x,y,z\in\ZZ[\rho]$ tals que $xyz\neq0$.
        \end{enumerate}
    \end{tcolorbox}
    \subsection{Exercici 3.1}
    Sigui $a,b\in \ZZ[\rho]$, amb $b\neq0$. Sigui $a/b=\alpha+\beta\rho$, i siguin $A$,$B$ els enters més propers a $\alpha$ i a $\beta$ respectivament. Sigui 
    \[
        N(x+iy)=(x+y\rho)(x+y(1-\rho))=x^2+xy+y^2
    \]
    la norma dels enters d'Eisenstein (la norma ve de $N(z)=z\sigma(z)$, on $\sigma$ envia $\rho\mapsto\rho^2=-1-\rho$). Aleshores, tenim que $N(\underbrace{A-\alpha}_{\eps_x})<\frac{1}{2}$, i $N(\underbrace{B-\beta}_{\eps_y})<\frac{1}{2}$. Per tant, tenim:
    \[
        a=\underbrace{(A+B\rho)}_qb+\underbrace{b(\eps_x+\eps_y\rho)}_r.
    \]
    Ara només ens falta veure que $N(r)<N(b)$:
    \[
        \frac{N(r)}{N(b)}=\frac{N(b(\eps_x+\eps_y\rho))}{N(b)}=N(\eps_x+\eps_y\rho)=
        \eps_x^2+\eps_x\eps_y+\eps_y^2\leq \frac{3}{4}<1
    \]
    Que és el que volíem veure.
    \subsection{Exercici 3.2}
    Notem que la norma de $\lambda$ és
    \[
    N(\lambda)=1+1+1=3
    \]
    Per tant, com que la norma és un primer de $\ZZ$, $\lambda$ ha de ser un primer de $\ZZ[\rho]$. Per veure que $3$ no és primer:
    \[
        \rho^{-1}\lambda^2=\frac{1-\sqrt 3i}{2}\left(\frac{3+\sqrt{3}i}{2}\right)^2=
        3
        % \frac{1-\sqrt 3i}{2}\frac{3+3\sqrt{3}i}{2}=
        % \frac{3(1-\sqrt 3i)(1+\sqrt{3}i)}{4}=
    \]
    Com que estem en un anell euclidià, tenim que és de factorització única, i per tant la definició de $\ord_\lambda$ té sentit. I l'argument que demostra que la valuació $p$-àdica satisfà aquesta propietat, també funciona aquí.
    \subsection{Exercici 3.3}
    Notem que aquest morfisme és lineal, per tant no ens hem de preocupar de que se satisfaci la propietat suma; només hem de veure que funciona amb el producte:
    \begin{align*}
        \Phi[(a+b\rho)(c+d\rho)]&=\Phi(ac+ad\rho+bc\rho+bd\rho-bd)\\
        &=ac-bd-ad-bc-bd\\
        &=a(c-d)-b(c-d)-3bd\\
        &=(a-b)(c-d)\\
        &=\Phi(a+b\rho)\Phi(c+d\rho)
    \end{align*}
    Que és el que volíem veure.\par
    Per veure la segona part, hem de veure que el nucli de $\Phi$ és $(\lambda)$ (i pel primer teorema d'isomorfia ja ho tenim); però:
    \begin{align*}
        \Phi(a+b\rho)=0&\iff a-b\equiv 0\pmod 3\\
        &\iff a=b+3k\\
        &\iff a+b\rho=b(1+\rho)+3k=b\lambda+\rho\inv\lambda^2k=\lambda(b+3k)\\
        &\iff a+b\rho\in(\lambda)
    \end{align*}
    Que és el que volíem veure.
    \subsection{Exercici 3.4}
    Suposem que $x^3+y^3=uz^3$, i que $\lambda\nmid x$, aleshores $x\equiv \pm 1\pmod \lambda$ (pel que hem vist a l'exercici anterior). Per tant:
    \[
        x^3=(\pm1+w\lambda)^3=\pm1 +3w\lambda\pm3w^2\lambda^2+w^3\lambda^3
    \]
    I mirat mòdul $\lambda^4$ obtenim:
    \[
        x^3\equiv\pm1 +\rho\inv w\lambda^3\pm \cancelto{0}{\rho\inv w^2\lambda^4}+w^3\lambda^3=\lambda^3w(1-\rho+w^2)\pm1=\lambda^3 w(2-\cancelto{0}{\lambda}+w^2)\pm1
    \]
    Que és el mateix que veure que $w(2+w^2)\equiv0\pmod \lambda$, però  els quadrats a $\ZZ_3$ són $0,1$, si és 0 ja està, i si és $1$ $2+w^2\equiv 0\pmod \lambda$, que és el que volíem veure.\par
    Ës a dir, si tenim que $\lambda\nmid xyz$, aleshores tindrem que $x,y,z\equiv \pm1\pmod \lambda$, i per tant, $x^3,y^3,z^3\equiv \pm1\pmod \lambda^4$, però no hi ha cap solució a $\pm1\pm1\equiv u\pmod \lambda^4$ (amb $u$ unitat), que és el que volíem veure.
    \subsection{Exercici 3.5}
    Suposem que $\lambda\nmid xy$, aleshores hi ha 3 casos: $x,y\equiv 1,1$, $x,y\equiv 1,-1$, $x,y\equiv -1,-1$. Si $x\equiv y$, aleshores ja sabem que això no té solució, i si $x\equiv-y\pmod \lambda$, tenim que $z^3\equiv 0\pmod\lambda^3$. Que això ens dona immediatament que $\lambda^2\mid z$.
    \subsection{Exercici 3.6}
    Suposem que $(x,y,z)$ és solució, aleshores $x,y\equiv 1,-1\pmod \lambda$.
    Considerem el primer cas: $\ord_\lambda(x+y)\geq1$ (no pot ser menor, ja que sinó no seria solució), aleshores sense pèrdua de generalitat podem escriure:
    \[
        x = 1+a\lambda\qquad y=-1+b\lambda
    \]
    Que ens dona que
    \[
        x\equiv 1+a\lambda \pmod {\lambda^2}\qquad
        y\equiv -1+b\lambda \pmod {\lambda^2}
    \]
    Però com que la suma ha de ser zero mòdul $\lambda^2$, tenim que $a\equiv-b\pmod \lambda^2$; i per tant:
    \[
        x\rho^2+y\rho^4\equiv (1+a\lambda)(\rho-1)+(1+a\lambda)\rho=(1+a\lambda)(2\rho-1)\not\equiv0\pmod{\lambda^4}.
    \]
    I l'altre cas es fa igual. Que ens dona el que volíem.\par
    Per fer el cas $\ord_\lambda(\rho^2x+\rho^4y)$, (o el cas $\ord_\lambda(\rho^4x+\rho^2y)$), notem que si enviem $x'=\rho^2x$ i $y'=\rho^4 y$, continuen sent solucions, per tant pel que hem vist abans, ja estem (notem que no pot ser que els tres tinguin ordre 1, ja que el producte té ordre almenys 4).
    \subsection{Exercici 3.7}
    Suposem que $(x,y,z)$ és solució primitiva, aleshores considerem
    \[
        \begin{array}{l}
            A=x+y\\
            B=\rho^2x+\rho^4y\\
            C=\rho^4x+\rho^2y\\
        \end{array}
    \]
    Aleshores, s'ha de veure que $A,B$ i $C$ són cubs (de fet, no ho són, però si dividim pel factor comú que tinguin, i en total tenim $ABC=uz^3$); i anomenem $z$ al element que té un $\ord_\lambda\geq4$, i als altres $x'$ i $y'$. Aleshores, és fàcil de veure que $x'+y'=z'$, i que $z'$ té un ordre menor que $z$. I fens un argument de descens infinit, arribes a que $x,y,z=0$.
    
    % \[
    %     \begin{cases*}
    %         x^3+y^3=uz^3\\
    %         (ax+b)^3+(cy+d)^3=uz/\lambda\\
    %         ax+b+cy+d = \frac{x+y}{\lambda}\\
    %         (ax+b)\rho^2+(cy+d)\rho^4 = \frac{x\rho^2+y\rho^4}{\lambda}\\
    %         (ax+b)\rho^4+(cy+d)\rho^2 = \frac{x\rho^4+y\rho^2}{\lambda}\\
    %         (\lambda-1)(ax+cy)=-\lambda(b+d)\\
    %         (\lambda-1)(ax\rho^2+cy\rho^4)=-\lambda(b\rho^2+d\rho^4) \\
    %         (\lambda-1)(ax\rho^4+cy\rho^2)=-\lambda(b\rho^4+d\rho^2) \\

    %     \end{cases*}
    % \]
    % \subsection{Exercici 3.8}
    % Sigui $(x,y,z)$ una solució primitiva de $x^3+y^3=uz^3$. Notem que fins ara sempre suposàvem que $\lambda\mid z$, però podria ser que $\lambda\mid x$:
    % \[
    %     (-x)^3+(-y)^3=u'z^3\iff x^3=y^3+(u''z)^3
    % \]
    % Que seria solució de l'equació de Fermat (a $\ZZ[\rho]$ hi han les arrels cúbiques de la unitat); per tant, sense pèrdua de generalitat, podem suposar que $$
    \newpage

% ------------------------------------------------------
% -------------------- Exercici 1 ---------------------
% ------------------------------------------------------
    \begin{tcolorbox}[title=\section{FP 2 Exercici 1}]
        Sigui $D$ un domini d'integritat i $x\in D\setminus \{0\}$. Demostreu que $x$ és irreductible si i només si $(x)$ és maximal entre tots els ideals principals propis de $D$. Deduïu que si $D$ és un domini d'ideals principals, aleshores $x$ és irreductible si i només si $(x)$ és maximal. Doneu un exemple d'un domini de factorització única $D$ i d'un element irreductible $x\in D$ tal que $(x)$ no sigui maximal.
    \end{tcolorbox}
    \subsection{Exercici 1}
    Notem que $x=ab$ si i només si $(x)\subseteq (a),(b)$. I d'aquí en traiem que si $x$ és irreductible ($a$ o $b$ són unitats), l'ideal és maximal; i viceversa. I en el cas de que $D$ sigui principal, aleshores ja estem.\par
    Considerem $\ZZ[X]$, notem que $X$ és irreductible, però $(X)\subsetneq (2,X)$.
    \newpage
% ------------------------------------------------------
% -------------------- Exercici 2 ---------------------
% ------------------------------------------------------
    \begin{tcolorbox}[title=\section{FP 2 Exercici 2}]
        Sigui $D=\ZZ[\frac{1+\sqrt{m}}{2}]$ per $m\equiv 1\pmod4$ i $m<0$.\par
        Sigui $D=\ZZ\left[\frac{1+\sqrt{m}}{2}\right]$ per $m\equiv 1\pmod4$ i $m<0$.
        \begin{enumerate}
            \item Demostreu que $u\in D^\times$ si i només si el valor absolut complex de $u$ és $1$.
            \item Demostreu que si $m<-3$ aleshores $D^\times=\pm1$.
            \item Demostreu que 2 i 3 són elements irreductibles de $\ZZ[\frac{1+\sqrt{-19}}{2}]$.
        \end{enumerate}
    \end{tcolorbox}
    \subsection{Exercici 2.1}
    Notem que $N(u)=u\ol u$, per tant, si té norma 1, $\ol u$ és l'invers. Si $u$ és unitat aleshores $N(uu\inv)=1=N(u)N(u\inv)$, i els 2 únics enters positius que multiplicats donen u són 1.\par
    \subsection{Exercici 2.2}
    Si $m<-3$, aleshores hem de resoldre:
    \[
        N(a+b\zeta)=(a+b\zeta)(a+b\ol \zeta)=1,
    \]
    on $\zeta=\frac{1+\sqrt{-m}}{2}$, per tant:
    \begin{align*}
        (a+b\zeta)(a+b\ol \zeta)=&\left(a+b\frac{1+\sqrt{-m}}{2}\right)\left(a+b \frac{1-\sqrt{-m}}{2}\right)\\
        =&\frac{(2a+b+b\sqrt{-m})(2a+b-b\sqrt{-m})}{4}\\
        =&\frac{4a^2+4ab+b^2-mb^2}{4}\\
        =&a^2+ab+\frac{1-m}{4}b^2=1
    \end{align*}
    I si mirem quins punts hi han a la bola $N\leq1$:
    \[
        \begin{cases*}
            b^2=\frac{4}{1-m}&\text{si }a=0\\
            1+\frac{1-m}4b^2+b=1&\text{si }a=1\\
        \end{cases*}
    \]
    Però si $m<-3$ aleshores $\frac{4}{1-m}<1$, i per tant $b^2<1$; que ens diu que només hi ha solució amb $a=1$ i $b=0$ per les $m<-3$.
    \subsection{Exercici 2.3}
    Notem que mirant normes no aconseguirem res. Per tant, considerem
    \[
        2=ab\implies N(2)=4=N(a)N(b)
    \]
    Per tant, hem de trobar elements de norma 2.
    \[
        2=N(x+y\zeta)=x^2+xy+10y^2
    \]
    Però $x^2+xy+y^2\geq0$, per tant tenim $9y^2<2$, és a dir, $y=0$, i per tant ha de ser real, però això no pot ser. Pel cas $3=ab$ tenim un cas similar, però amb $9y^2<3$.
    \newpage
% ------------------------------------------------------
% -------------------- Exercici 3 ---------------------
% ------------------------------------------------------
    \begin{tcolorbox}[title=\section{FP 2 Exercici 3}]
        Sigui $D$ un domini d'integritat que no és un cos. Diem que un element $\beta\in D\setminus(D^\times\cup\{0\})$ és un \textit{divisor lateral universal} si per a tot $\alpha\in D$ existeix $\delta\in D^\times\cup\{0\}$ tal que $\alpha-\delta$ és divisible per $\beta$.
        \begin{enumerate}
            \item Demostreu que si $D$ no té divisors laterals universals, aleshores $D$ no és domini euclidià respecte a cap funció euclidiana
            \[
                N:D\to\ZZ_{\geq0}.
            \]
            \textit{Indicació: Suposant que $D$ és un domini euclidià respecte a $N$, considereu el mínim del conjunt \[S=\{N(\alpha)\mid\alpha\in D\setminus(D^\times\cup\{0\})\}.\]}
            \item Demostreu que $D=\ZZ\left[\frac{1+\sqrt{-19}}{2}\right]$ no és domini euclidià per a cap funció euclidiana.
            \textit{Indicació: Considerant l'element $\alpha=2$, vegeu que els únics divisors laterals universals possibles són associats a 2 i 3. Considerant l'element $\zeta=\frac{1+\sqrt{-19}}{2}$, demostreu que 2 i 3 no poden ser divisors laterals universals.}\par
        \end{enumerate}
        \textit{Comentari: En un full de problemes posterior demostrarem, no obstant, que $D$ és un domini d'ideals principals.} 
    \end{tcolorbox}
    \subsection{Exercici 3.1}
    Suposem que $D$ és un domini euclidià. Sigui $x$ un element de norma mínima de $D\setminus(D^\times\cup\{0\})$. Considerem ara un $y\in D$ qualsevol, com que $D$ és euclidià, sabem que existeixen $q,r\in D$ tals que $y=xq+r$ amb $N(r)<N(x)$, i com que $x$ té norma mínima, $r\in D^\times\cup\{0\}$.
    \subsection{Exercici 3.2}
    Suposem que $D$ té un divisor lateral universal $z$, per tant:
    \[
        z\mid 2+u\qquad \text{on $u$ és una unitat.}
    \]
    És a dir:
    \[
        z\mid 1\quad\text{ o }\quad z\mid 2\quad \text{ o }\quad z\mid3.
    \]
    És a dir, $z\in \{\pm2,\pm3\}$ (i per tant, la norma pot ser 4 o 9). Com que ha de ser un divisor lateral per tot $d\in D$, en particular per $d=\zeta$:
    \[
        z\mid \frac{1+\sqrt{-19}}{2}
        \quad
        \text{ o }
        \quad
        z\mid \frac{1+\sqrt{-19}}{2}+1
        \quad
        \text{ o }
        \quad
        z\mid \frac{1+\sqrt{-19}}{2}-1.
    \]
    I mirant normes, tenim:
    \[
        N(z)\mid 5
        \quad
        \text{ o }
        \quad
        N(z)\mid 7.
    \]
    Però que això contradiu el que dèiem abans.
    \newpage
% ------------------------------------------------------
% -------------------- Exercici 1 ---------------------
% ------------------------------------------------------
    \begin{tcolorbox}[title=\section{FP 4 Exercici 1}]
        Sigui $K=\QQ(\sqrt{m})$, on $m$ és un enter lliure de quadrats. Demostreu que l'anell d'enters algebraics de $K$ està donat per
        \[
            \mathcal{O}_K=\begin{cases}
                \ZZ[\sqrt m]& \text{si }m\not \equiv 1\pmod4,\\
                \ZZ[\frac{1+\sqrt m}2]& \text{si }m \equiv 1\pmod4,\\
            \end{cases}
        \]
    \end{tcolorbox}
    \subsection{Exercici 1}
    Primer veurem la primera inclusió ($\supseteq$):\\
    Notem que $\sqrt{m}$ és arrel de $x^2-m$. I si $m\equiv1\pmod4$, aleshores $\half(1+\sqrt{m})$ és arrel de $x^2-x-\frac{m-1}{4}$.\par

    Per veure l'altra inclusió, prenem $a\in \QQ(\sqrt{m})$ un enter algebraic. Aleshores $a=x+y\sqrt{m}$ amb $x,y\in \QQ$, i com que és enter, existeix un polinomi $f$ que l'annul·la. Com que $[\QQ(\sqrt{m}):\QQ]=2$, podem suposar que $f$ és un polinomi mònic de grau dos: $f(X)=X^2+pX+q$, amb $p,q\in \ZZ$. Per tant:
    % \[
    %     f(a)=
    %     % (x+y\sqrt{m})^2+p(x+y\sqrt{m})+q=
    %     % x^2+2xy\sqrt{m}+my^2+px+py\sqrt{m}+q=
    %     (x^2+my^2+px+q)+(2x+p)y\sqrt{m}=0
    % \]
    % I per tant:
    % \[
    %     \begin{cases}
    %         x^2+px+q+my^2=0\\
    %         (2x+p)y=0
    %     \end{cases}
    % \]
    % Notem que si $y=0$, aleshores $a\in \ZZ$; per tant podem suposar que $y\neq0$.
    % \[
    % \begin{cases}
    %     x^2+px+q+my^2=0\\
    %     x=\frac{-p}{2}
    % \end{cases}\implies
    % \begin{cases}
    %     \frac{4q-p^2}{4}+my^2=0\\
    %     x=\frac{-p}{2}
    % \end{cases}\implies
    % \]
    \[
        2a=-p\pm\sqrt{p^2-4q}=2x+2y\sqrt{m}
    \]
    Si suposem que $a\notin\ZZ$, aleshores tenim que
    \[
        \begin{cases}
            -p=2x\\
            p^2-4q=4y^2m
        \end{cases}
    \]
    De l'igualtat de dalt, en podem treure que $x\in\half\ZZ$. De l'igualtat de baix en podem treure que $4y^2m\in \ZZ$. Però si escrivim $y=\frac{y_n}{y_d}$, aleshores, $y_d^2\mid 4m$, però com que $m$ és lliure de quadrats, tenim que $y_d\mid 2$: és a dir, $y\in \half\ZZ$.\par
    Sigui $2x=u\in \ZZ$ i $2y=v\in\ZZ$. Aleshores:
    \[
        u^2-mv^2=4x^2-4y^2m=4q\in4\ZZ
    \]
    Per tant, si ens mirem les coses mòdul 4 obtenim que
    \[
    u^2\equiv mv^2\pmod4
    \]
    Si $u^2\equiv 0\pmod 4$, aleshores, com que $4\nmid m$, $v^2\equiv 0\pmod 4$ (i per tant, els dos són parells, per tant $x+y\sqrt{m}\in \ZZ[\sqrt{m}]$).\\
    Si $u^2\equiv 1\pmod 4$, aleshores $v^2\equiv 1\pmod 4$ i $m\equiv1\pmod4$. I d'aquí podem concloure que
    \[
        x+y\sqrt{m}=\frac{u+v\sqrt{m}}{2}=\frac{u-v}{2}+v\frac{1+\sqrt{m}}{2}
    \]
    (Com que $u\equiv v\pmod 2$, la resta es pot dividir per 2).
    \newpage
% ------------------------------------------------------
% -------------------- Exercici 2 ---------------------
% ------------------------------------------------------
    \begin{tcolorbox}[title=\section{FP 4 Exercici 2}]
        Sigui $K=\QQ(\sqrt{m})$, on $m$ és un enter \emph{negatiu} lliure de quadrats. Demostreu que
        \[
            (\mathcal{O}_K)^\times=\begin{cases}
                \langle i \rangle&\text{si } m=-1,\\
                \langle \rho \rangle&\text{si } m=-3,\\
                \langle -1 \rangle&\text{altrament.}
            \end{cases}
        \]
    \end{tcolorbox}
    \subsection{Exercici 2}
    Notem que si $m\not\equiv1\pmod 4$, aleshores 
    \[
        N(a)=x^2-my^2\geq x^2+y^2
    \]
    Per tant, només hem de comprovar $\pm1$ (que són unitats), i $\pm \sqrt{m}$ que és unitat si i només si $m=-1$.\par
    El cas $m\equiv 1\pmod 4$ ja s'ha fet en una altra llista de problemes.
    \newpage
% ------------------------------------------------------
% -------------------- Exercici 3 ---------------------
% ------------------------------------------------------
    \begin{tcolorbox}[title=\section{FP 4 Exercici 3}]
        Sigui $K=\QQ(\sqrt{m})$, on $m\neq1$ és un enter \emph{positiu} lliure de quadrats. L'objectiu d'aquest exercici és demostrar que
        \[
            \#(\mathcal{O}_k)^\times=\infty
        \]
        Sobre $\ZZ[\sqrt{m}]$, definim la següent norma
        \[
            N:\ZZ[\sqrt{m}]\to\ZZ,\qquad N(a+b\sqrt{m})=a^2-mb^2.
        \]
        Noteu que aquest valor no és el quadrat de la norma complexa i que, en particular, pot prendre valors negatius.
        \begin{enumerate}
            \item Sigui $\alpha\in\ZZ[\sqrt{m}]$. Demostreu que $\alpha\in\ZZ[\sqrt{m}]^\times$ si i només si $N(\alpha)=\pm1$. Doneu un exemple de $m$ i $\alpha$ tal que $N(\alpha)=-1$.
            \item Demostreu que per a tot $n\in\ZZ_{n\geq1}$ existeixen $a,b\in\ZZ$ tals que\footnote{Aquesta desigualtat es pot fer més estricta, canviant el $1/n$ per $1/n^2$, pel teorema d'aproximació de Dirichlet.}
            \begin{equation}\label{FP4P3 aproximacio de arrels quadrades}
                0<|a-b\sqrt{m}|<\frac{1}{n}\quad\text{ i }\quad 0<b<n.
            \end{equation}
            \item Demostreu que existeixen infinits parells d'enters $(a,b)$ tals que
            \begin{equation}\label{FP4P3 bones aproximacions tenen norma petita}
                0<|N(a+b\sqrt{m})|<1+2\sqrt{m}.
            \end{equation}
            \item Demostreu que existeix un enter $t$ amb $0<|t|<1+2\sqrt{m}$ tal que l'equació
            \[
                N(a+b\sqrt{s})=t
            \]
            té infinites solucions en $a,b\in\ZZ$ amb $a,b>0$ i $\gcd(b,t)=1$.
            \item Demostreu que existeix $\alpha \in\ZZ[\sqrt{m}]$ amb $\alpha\neq \pm1$, tal que $N(\alpha)=\pm1$.
            \item Deduïu de l'apartat anterior que $\#(\mathcal{O}_K)^\times=\infty$.
        \end{enumerate}
    \end{tcolorbox}
    \subsection{Exercici 3.1}
    Per veure la primera implicació, si $N(a)=a\ol a=1$.\par
    Si $ab=1$ aleshores $N(a)N(b)=1$ i per tant $N(a)\in\ZZ^\times=\langle-1\rangle$.\par
    A $\ZZ[\sqrt{2}]$, $\alpha=1-\sqrt{2}$ té norma $-1$.
    \subsection{Exercici 3.2}
    Considerem el conjunt 
    \[A_n=\left\{\floor{i\sqrt{m}}-i\sqrt{m}+1\mid i=0,1,\dots,n\right\}.\]
    Notem que $|A_n|=n+1$, i que
    \[
        0<\floor{i\sqrt{m}}-i\sqrt{m}+1\leq1
    \]
    Per tant, si trenquem l'interval $(0,1]$ en $n$ trossos ($(0,1]=(0,1/n]\cup(1/n,2/n]\cup\dots\cup(n-1/n,1]$) hi haurà com a mínim dos elements de $A_n$ en el mateix interval. Suposem que $\floor{i\sqrt{m}}-i\sqrt{m}+1$ i $\floor{j\sqrt{m}}-j\sqrt{m}+1$ pertanyen al mateix interval, aleshores
    \[
        \underbrace{\floor{i\sqrt{m}}-\floor{j\sqrt{m}}}_{a}
        -\underbrace{(i-j)}_b\sqrt{m}.
    \]
    Per veure que $0<b<n$, notem que ara mateix només tenim que $-n<b<n$ i que $b\neq0$, per tant, si $b<0$ podem prendre $a\mapsto -a$ i $b\mapsto-b$. Notem que hem pres $a$ i $b$ de tal manera que se satisfà la propietat de l'enunciat.
    \subsection{Exercici 3.3}
    Fixem un $n>0$, i suposem que tenim un $(a,b)$ que satisfà la propietat \ref{FP4P3 aproximacio de arrels quadrades}, volem veure que aleshores satisfà també la propietat \ref{FP4P3 bones aproximacions tenen norma petita} (com que $m$ és lliure de quadrats, aleshores $N(p)=p\ol p$).\par
    \[
        |N(a+b\sqrt{m})|=
        |a+b\sqrt{m}||a-b\sqrt{m}|<
        \frac{|a+b\sqrt{m}|}{n}=
        \left|\frac{a}{n}+\frac{b}{n}\sqrt{m}\right|
    \]
    Però del fet que $0<b<n$ en podem deduir que:
    \begin{align}\label{FP4P3 desigualtat 1}
        |a/n+b/n\sqrt{m}|\leq|a/n|+|b/n\sqrt{m}|\leq|a/n|+\sqrt{m}.
    \end{align}
    I de que $|a-b\sqrt{m}|<1/n$ en treiem que
    \begin{align}\label{FP4P3 desigualtat 2}
        |a/n|\leq\frac{|a-b\sqrt{n}|+|b\sqrt{m}|}{n}\leq 1/n^2+\sqrt{m}<1+\sqrt{m}.
    \end{align}
    I ajuntant \ref{FP4P3 desigualtat 1} i \ref{FP4P3 desigualtat 2} obtenim que $N(a+b\sqrt{m})<1+2\sqrt{m}$, i com que $b\sqrt{m}\not\in \ZZ$ la norma no pot ser zero.\par
    Per tant, si fem variar la $n$ cada vegada més petita, obtindrem una successió amb termes diferents. És a dir, si considerem $a_k+b_k\sqrt{m}$ un element que satisfaci \ref{FP4P3 aproximacio de arrels quadrades} per $n=k$, aleshores podem considerar $\{a_{k_j}+b_{k_j}\sqrt{m}\}_{j\in\NN}$ una successió tal que $\frac{1}{k_{j+1}}<|a_{k_j}-b_{k_j}\sqrt{m}|(<\frac{1}{k_j})$, aleshores tots els termes seran diferents, i per tant, sabem que n'hi ha infinits.
    \subsection{Exercici 3.4}
    Considerem la seqüència $a_n+b_n\sqrt{m}$ com l'hem definit abans. Notem que podem suposar que $a_n$ i $b_n$ són coprimers (si no ho són, dividint aconseguim una nova parella amb norma més petita, i amb $b'<b$). Prenent la parcial $a_{n_k}+b_{n_k}\sqrt{m}$ com l'hem definit abans.\par
    Com que tenim infinits termes, i la norma és un enter, per tant hi ha d'haver algun enter $1,2,\dots,\floor{1+2\sqrt{m}}$ que es repeteixi infinites vegades. Ara hem de veure que $a,b>0$ i que $\gcd(b,t)=1$. Notem que podem canviar el signe de $a$ i de $b$ sense que canviï la norma, per tant, podem suposar que $a,b>0$.\par
    Sigui $g=\gcd(b,t)$, aleshores
    \[
        g=\gcd(b,t)=\gcd(b,a^2-b^2m)=\gcd(b,a^2)\leq\gcd(b,a)^2=1
    \]
    Per tant, si $g=1$, que és el que volíem veure.
    \subsection{Exercici 3.5}
    Considerem la successió $\{(a_n,b_n)\}_n$ solucions de l'equació anterior. Considerem $a_ib_i^{-1}\pmod t$, una d'aquestes classes tindrà infinits elements. Sigui $\{(a_{n_k},b_{n_k})\}_k$ una successió (infinita) d'elements tals que $a_{n_i}b_{n_i}\inv\equiv a_{n_1}b_{n_1}\inv\pmod t$. Sigui $(a_{n_i},b_{n_i})=(a,b)$ i $(a_{n_j},b_{n_j})=(a',b')$. Considerem
    \[
        \alpha=\frac{a+b\sqrt{m}}{a'+b'\sqrt{m}}=\frac{(a+b\sqrt{m})(a'-b'\sqrt{m})}{t}=\frac{aa'-bb'm}{t}+\frac{a'b-ab'}{t}\sqrt{m}
    \]
    Però notem que $t\mid a'b-b'a$. De $N(a+b\sqrt{n})=t$ i de $ab'\equiv a'b\pmod t$ en podem treure:
    \[
        aa'\equiv ab' \frac{a'}{b'}\equiv a^2\frac{b'}{b}\equiv b^2m\frac{b'}{b}=bb'm 
    \]
    Per tant, $\alpha\in\mathcal{O}_k$; i per simetria, $\alpha\inv\in\mathcal{O}_k$, que ens diu que $\alpha\in\mathcal{O}_k^\times$.
    \subsection{Exercici 3.6}
    Per la construcció d'abans, tenim que per cada parella de $(a_{n_k},b_{n_k})$, $\alpha_{(i,1)}=\frac{a_{n_i}+b_{n_i}\sqrt{m}}{a_{n_1}+b_{n_1}\sqrt{m}}$. Per tant, tenim que $\mathcal{O}_K^\times$ és un conjunt infinit.
    \newpage
% ------------------------------------------------------
% -------------------- Exercici 4 ---------------------
% ------------------------------------------------------
    \begin{tcolorbox}[title=\section{FP 4 Exercici 4}]
        Escriviu
        \[
            \alpha=\left(\frac{1+\sqrt{2}}{9}\right)^{1/3}+\left(\frac{1-\sqrt{2}}{9}\right)^{1/3}
        \]
        com el quocient d'un enter algebraic per un enter.
    \end{tcolorbox}
    \subsection{Exercici 4}
    Considerem $\alpha=a/3$. Aleshores, volem veure que $a$ és un enter algebraic. Recordem que els enters algebraic són un anell, per tant, si veiem que els dos sumands són algebraics, ja estarem.\par
    Recordem que que a teoria vam veure que si un nombre és arrel d'un polinomi mònic amb coeficients a $\ol\ZZ$, aleshores és un enter algebraic.\par
    Per veure que $\sqrt{2}$ és un enter algebraic, considerem $T^2-2$. Per veure que $(3\pm3\sqrt{2})^{1/3}$ és un enter algebraic considerem $T^3-3\mp3\sqrt{2}$. Que és el que volíem veure.
    % \begin{align*}
    %     \alpha&=\left(\frac{1+\sqrt{2}}{9}\right)^{1/3}+\left(\frac{1-\sqrt{2}}{9}\right)^{1/3}=\\
    %     &=\left(
    %         \frac{2}{9}
    %         +\frac{1}{3}\sqrt[3]{
    %         -1-\sqrt{2}}
    %         +\frac{1}{3}\sqrt[3]{-1+\sqrt{2}}
    %     \right)^{1/3}.
    % \end{align*}
    % Per tant $3\alpha$ es pot escriure com 
    % \[
    %     3\alpha=\left(6+9\sqrt[3]{-1-\sqrt{2}}+9\sqrt[3]{-1+\sqrt{2}}\right)^{1/3}
    % \]
    % Ara veurem que $3\alpha$ és algebraic. Recordem que a teoria vam veure que si un nombre és arrel d'un polinomi mònic amb coeficients enters algebraics, aleshores és un enter algebraic. Per tant, considerem el polinomi:
    % \[
    %     \begin{array}{r@{\,=\,}l}
    %         p_1(X)&X^2-2\\
    %         p_2(X)&X^3+1+\sqrt{2}\\
    %         p_3(X)&X^3+1-\sqrt{2}\\
    %         p_4(X)&X^3-6-9\sqrt[3]{-1-\sqrt{2}}-9\sqrt[3]{-1+\sqrt{2}}
    %     \end{array}
    % \]
    % I com que $3\alpha$ és arrel de $p_4$
    % \[
    %     \frac{x^9 - 18x^6 + 108x^3 + 1242}{3^{21}}=
        
    %     \sqrt[3]{1-\sqrt{2}}
    %     +\sqrt[3]{1+\sqrt{2}}
    % \]
    % \[
    %     % (((X^3-6)/9)^3+2)/3=
    %     \frac{(x^9 - 18x^6 + 108x^3 + 1242)^3}{3^{21}}=
    %     2+\sqrt[3]{-1+\sqrt{2}}+\sqrt[3]{-1-\sqrt{2}}
    % \]
    \newpage
\end{document}
% ------------------------------------------------------
% -------------------- Exercici XX ---------------------
% ------------------------------------------------------
    % \begin{tcolorbox}[title=\section{FP X Exercici XX}]
    % \end{tcolorbox}
    % \subsection{Exercici XXX}
    % \newpage
